\documentclass[12pt]{article}

%% Basic document formatting
\usepackage{amsmath, amsthm, amssymb, amsfonts}
\usepackage{mathtools}
\usepackage{mathrsfs}
\usepackage{xspace}
\usepackage{thmtools}
\usepackage{graphicx}
\usepackage{tikz-cd}
\usepackage{color}
\usepackage{setspace}
\usepackage{fancyhdr}
\usepackage{titling}
\usepackage{hyperref}
\usepackage[left=0.4in,right=0.4in,top=1in,bottom=1in]{geometry}
\usepackage{float}
\usepackage{tabularx}
\usepackage{booktabs}
\usepackage[utf8]{inputenc}
\usepackage[english]{babel}
\usepackage{framed}
\usepackage[dvipsnames]{xcolor}
\usepackage{environ}
\usepackage{tcolorbox}
\tcbuselibrary{theorems,skins,breakable}

\usepackage{enumitem}
\setlist[enumerate]{leftmargin=*}
\setlist[enumerate,1]{labelindent=\parindent}
\setlist[enumerate,2]{labelindent=0pt}

% Blackboard Bold
\newcommand{\Z}{\mathbb{Z}}                 % integers
\newcommand{\Zpl}{\mathbb{Z}_{+}}           % positive integers
\newcommand{\Znn}{\mathbb{Z}_{\geq 0}}      % nonnegative integers. 
\newcommand{\Q}{\mathbb{Q}}                 % rationals
\newcommand{\Qpl}{\mathbb{Q}_{+}}           % positive Rationals
\newcommand{\R}{\mathbb{R}}                 % reals
\newcommand{\Rpl}{\mathbb{R}_{+}}           % positive Reals
\newcommand{\C}{\mathbb{C}}                 % complex numbers
\newcommand{\F}{\mathbb{F}}                 % field

% Words
\newcommand{\st}{\text{ such that }}
\newcommand{\wrt}{\text{ with respect to }}
\newcommand{\with}{\text{ with }}
\newcommand{\ie}{\text{, i.e. }}
\newcommand*{\ditto}{\text{---}\texttt{"}\text{---}}

% Operators
\newcommand{\abs}[1]{\left|#1\right|}                   % absolute value
\newcommand{\norm}[1]{\abs{\abs{#1}}}
\newcommand{\floor}[1]{\left\lfloor #1 \right\rfloor}   % floor
\newcommand{\ceil}[1]{\left\lceil #1 \right\rceil}      % ceiling
\newcommand{\powset}[1]{\mathscr{P}(#1)}

% Category Theory 
\renewcommand{\hom}{\text{Hom}}
\newcommand{\homspace}[3]{\hom_{#1}(#2, #3)}
\newcommand{\coproduct}[9]{
  \begin{tikzcd}[ampersand replacement=\&]
      \& \& #1 \arrow[dll, bend right, "#6"] \arrow[dl, "#5"] \\
   #4 \& #3 \arrow[l, "#9"] \\
      \& \& #2 \arrow[ull, bend left, "#8"] \arrow[ul, "#7"]
  \end{tikzcd}
}

\newcommand{\product}[9]{ 
  \begin{tikzcd}[ampersand replacement=\&]
      \& \& #3 \\
      #1 \arrow[r, "#7"] \arrow[urr, bend left, "#5"] \arrow[drr, bend right, "#6"] \& #2 \arrow[ur, "#8"] \arrow[dr, "#9"] \\ 
      \& \& #4
  \end{tikzcd}
}


% Algebra 
\newcommand{\Syl}[2]{\text{Syl}_{#1}{(#2)}}           % set of Sylow-p subgroups 
\newcommand{\<}{\langle}                % \<x\>, subgroup generated by x 
\renewcommand{\>}{\rangle}
\renewcommand{\index}[2]{[#1 : #2]}
\newcommand{\id}{\text{id}}             % identity element
\newcommand{\lhdneq}{%
  \mathrel{\ooalign{$\lneq$\cr\raise.22ex\hbox{$\lhd$}\cr}}}
\newcommand{\order}[1]{\text{o}(#1)}    % order of an element
\newcommand{\spec}{\operatorname{Spec}}
\newcommand{\rad}{\operatorname{rad}}   % radical of an ideal 
\newcommand{\sym}[1]{\mathfrak{S}_{#1}}
\newcommand{\aut}[1]{\text{Aut}(#1)} 
\newcommand{\gal}[2]{\text{Gal}(#1 / #2)} 
\newcommand{\inn}[1]{\text{Inn}(#1)} 
\let\oldcong\cong
\let\oldequiv\equiv
\renewcommand{\cong}{\oldequiv}
\renewcommand{\equiv}{\oldcong}

\newcommand{\triangleleftneq}{\mathrel{\ooalign{%
  $\trianglelefteq$\cr
  \hidewidth\raise0.06ex\hbox{$\not\mkern4mu$}\hidewidth\cr
}}}

\makeatletter % cycle
\newcommand{\cyc}[1]{(\mathbf{\cyc@process#1\relax})}
\def\cyc@process#1#2\relax{%
	#1%
	\ifx\relax#2\relax
	\else
		\,\cyc@process#2\relax
	\fi
}
\makeatother

\newcommand{\GL}[2]{\text{GL}_{#1}(#2)}                             % general linear group
\newcommand{\SL}[2]{\text{SL}_{#1}(#2)}                             % special linear group
\newcommand{\gl}[2]{\mathfrak{gl}_{#1}(#2)}
\renewcommand{\sl}[2]{\mathfrak{sl}_{#1}(#2)}
\newcommand{\so}[2]{\mathfrak{so}_{#1}(#2)}
\newcommand{\su}[1]{\mathfrak{su}(#1)}

% Linear Algebra
\newcommand{\bmat}[1]{\begin{bmatrix}#1\end{bmatrix}}   % bracketed matrix
\newcommand{\rank}{\operatorname{rank}}                 % rank
\newcommand{\nullity}{\operatorname{nullity}}           % nullity
\newcommand{\tr}{\operatorname{tr}}

% Combinatorics
\newcommand{\qnum}[1]{\left[#1\right]_q}                % q-analog
\newcommand{\qfact}[1]{\left[#1\right]!_q}              % q-analog factorial 
\newcommand{\qbinom}[2]{\binom{#1}{#2}_q}               % Gaussian binomial coefficient

% Topology/Analysis
\newcommand{\ball}[2]{\text{B}_{#1}(#2)}  % B_r(x): open r-balls around x
\newcommand{\diam}{\text{diam}}           % diamter of a set in metric space 
\newcommand{\lowint}[2]{\underline{\int_{#1}^{#2}}}
\newcommand{\upint}[2]{\overline{\int}_{#1}^{#2}}

% Misc Notation
\newcommand{\oneton}[1]{\{1, \ldots, #1\}}
\newcommand{\defeq}{\vcentcolon=}   % :=
\newcommand{\eqdef}{=\vcentcolon}   % =: 
\renewcommand{\bf}[1]{\textbf{#1}}
\newcommand{\im}{\operatorname{im}}
\newcommand{\fnrest}[2]{\left.#1\right|_{#2}}
\newcommand{\isom}{\overset{\sim}{\rightarrow}} 


\renewcommand{\thesection}{\arabic{section}.}
\renewcommand{\thesubsection}{(\alph{subsection})}

\newtheorem{claim}{Claim}
\newtheorem*{lemma}{Lemma}

%% Headers & title setup
\newcommand{\course}{Math100C}
\newcommand{\myname}{Jay Ser}
\setlength{\headheight}{14.5pt}
\pagestyle{fancy}
\fancyhf{}
\renewcommand{\headrulewidth}{0.4pt}
\lhead{\course}
\rhead{\myname}
\cfoot{\thepage}
\setlength{\droptitle}{-4em} 
\title{\course\ - WK7} 
\author{\myname}
\date{2026.05.17}

\begin{document}
\maketitle
\thispagestyle{fancy}

\newtheorem*{hint}{Hint}

%------------------------------------------------------------------------------%

\section{} 
Take field extensions $\F_{t} / \F_q$ of finite fields. 
These fields must have the same character, namely of some prime $p$. 
That is, one really has field extensions 
$$\F_{t} / \F_{q} / \F_{p}$$ 
It was shown in week 6 that the primitive element theorem holds for finite extensions of $\F_p$, so there is some $\alpha \in \F_{t}$ such that $\F_{t} = \F_{p}(\alpha)$. 
Then 
$$\F_{t} = \F_{q}(\alpha);$$
By definition, $\F_{q}(\alpha) \subset \F_{t}$. 
To get containment in the other direction, $\F_{q}(\alpha)$ contains $\F_p$ and $\alpha$, i.e., 
$\F_{t} = \F_{p} (\alpha) \subset \F_{q}(\alpha)$. 

\section{} 
The backward direction follows because the Frobenius map is an automorphism of $F$: 
$$c(u + v)^{p^n} = cu^{p^n} + cv^{p^n},$$ 
hence if each term of $f$ has degree $c_i x^{p^{n_i}}$, then $u \mapsto f(u)$ is a homomorphism on $F$. 

\section{} 
Recall the fundamental theorem for finitely generated modules over a PID $R$: 
if $M$ is such a module, there exist a number $r, m$ and nonunit elements $a_1 \mid a_2 \mid \ldots \mid a_m$ such that 
$$M \equiv R^{\oplus r} \oplus R / \<a_1\> \oplus R / \<a_2\> \oplus \cdots \oplus R / \<a_m\>$$
By definition, finitely generated free modules over $R$ are those of the form $R^{\oplus r}$. 
Suppose $M$ is finitely generated. 
If $M$ is free, then because $R$ is a PID and thus particularly an integral domain, $M = R^{\oplus r}$ has no torsion elements. 

Vice versa, if $M$ is torsion-free, then all the $R / \<a_i\>$ factors need to be the zero module, for $R / \<a_i\>$ as a $R$-module has torsion elements:
namely, for any $b + \<a_i\> \in R / \<a_i\>$, $a_i \cdot (b + \<a_i\> = 0$ in $R / \<a_i\>$. 
So $M \equiv R^{\oplus r}$, which proves $M$ is a free $R$-module. 

If $R$ is any ring, recall that considering $R$ as a module over itself, the submodules of $R$ are precisely the ideals in $R$. 
Let $R = \C[x, y]$. 
Consider the ideal $\<x, y\> \lhd \C[x]$ as an $R$-module. 
It is clearly torsion-free because $R$ is an integral domain. 
It is also inuitively clear that it is a free-module, though I can't find a proof... 

\section{}
I saw this proof, which is also the proof intended by Kiran, while going through Dummit and Foote. 
Suppose $\Phi_{n}(x) = f(x)g(x)$ is a factorization of $\Phi_n$, where $f$ is irreducible over $\Q$. 
Fix a primitive $n$th root of unity $\zeta \in \C$ that is a root of $f$ (i.e., $f$ is the minimal polynomial for $\zeta$ over $\Q$; 
all roots of $f$ are a primitive $n$th root of unity because they are, by construction, the roots of $\Phi_n(x)$), and take any prime $p$ not dividing $n$. 
Suppose $\zeta^p$ is a root of $g$. 
Then $\zeta$ is a root of $g(x^p)$, so $g(x^p) = f(x) h(x)$ for some $h \in \Z[x]$ (by Gauss Lemma). 
Over $\F_p$, $\overline{g}(x^p) = \overline{g}(x)$ gives 
$$\overline{g}(x) = \overline{f}(x) \overline{h}(x)$$
Hence over $\F_p$, $\overline{\Phi}_n(x)$ has a repeated factor; 
since $\Phi_n(x)$ is a factor of $x^n - 1$ in $\Z[x]$, it follows $x^n - 1$ has a repeated factor over $\F_p$. 
This contradicts the obvious observation that $x^n - 1$ doesn't have repeated roots over $\F_p$ since its derivative is $nx^{n - 1}$, which isn't 0 since $p \nmid n$; 
$x^n - 1$ and its derivative have no common factor, so $x^n - 1$ can't have a repeated root. 
This shows that $\zeta^p$ is not a root of $g$; since $\zeta^p$ is a root of $\Phi_n(x)$, it follows $\zeta^p$ is a root of $f$. 

Let $a$ be coprime to $n$. 
Then $a = p_1 \cdots p_k$ for primes $p_i \nmid n$. 
Since the argument in the paragraph above was for an arbitrary root $\zeta$ of $f$ and any prime not dividing $n$, $\zeta^{a}$ is a root of $f$; 
the conclusion is 
$$\zeta^{a} \text{ is a root of } f \text{ for every } a \text{ coprime to } n$$
i.e., every primitive $n$th root of unity is a root of $f$. 
This forces $f = \Phi_n$; $\Phi_n(x)$ is irreducible. 



\end{document}
